\section{Starting point}
\label{sec:start}

\subsection{What already exists}

\begin{center}
\small
\begin{tabularx}{\textwidth}{@{}L{5.0cm}Y@{}}
\toprule
\textbf{Asset (location)} & \textbf{State and how A2 uses it} \\
\midrule
\fp{asm1_demo/explorer/esplora.py} & A polite Esplora client using the standard library only. It pages with the last-seen-txid cursor, waits between requests, backs off on HTTP 429 and 5xx, and keeps request statistics. Reused unchanged for DS1 and DS2. \\
\fp{asm1_demo/explorer/amounts.py} and \fp{tests/test_amounts.py} & Received, sent and net per address, with 9 passing tests. The D1 check showed that summing them reproduces the API's own \code{funded_txo_sum} and \code{spent_txo_sum} (risk 18). Reused unchanged. \\
\fp{asm1_demo/config/settings.json} & Every threshold in one file: selection rule, windows, minimum history, MAD scale, flag at 3.5, fit fraction 0.7, top 2\%, sanity 20\%, banned words. It becomes the A2 configuration and gains new sections for the models. \\
\fp{asm1_demo/frontend/src/api/mock/scoring.ts} & A complete reference implementation of F1--F4 and the baseline, including CR-02, CR-04 and CR-05 and the same-block rule. The Python \code{features.py} is a line-by-line port, so the prototype and the pipeline agree. \\
\fp{asm1_demo/scripts/d1_check.py}, \fp{reports/d1-check.md} & One address (59 transactions) checked end to end. Its seeded block-sampling logic is the starting point for the DS1 selection script. \\
\fp{asm1_demo/frontend/src/api/types.ts} & The JSON shapes the A3 front end expects. \code{predict.py} returns the same shapes. \\
\fp{asm2/lm/} (model, tokenizer, data, train, chat, generate, smoke) & A GPT with RoPE, RMSNorm and SwiGLU, its own BPE tokenizer, memory-mapped data, masked chat fine-tuning, a speed benchmark, and Ctrl+C save and resume. It passes 15 tests and an 8-stage smoke test on CPU. It has not yet run on the GPU. The same model class serves TxGPT. \\
\fp{asm2/configs/text_150m_pretrain.json}, \fp{text_150m_chat.json} & 152.6M parameters (18 layers, $d=768$, 12 heads, sequence 1,024). 2B pretraining tokens, then a 100M-token chat fine-tune. \\
\bottomrule
\end{tabularx}
\end{center}

\subsection{What is missing}

\begin{enumerate}
  \item \textbf{The real snapshot.} No folder contains the 25 + 3 address snapshot. The A1 Gantt
    chart planned WBS 1.2.1 for 11--13 October, so the critical path is about one week behind.
    The work itself is short (Section~\ref{sec:ds1}: under two hours of fetching), so this is
    recoverable if it starts on 21 October.
  \item \textbf{The whole scikit-learn pipeline.} There are no features, models, evaluation,
    notebooks or figures yet.
  \item \textbf{Model types that match the brief.} Isolation Forest and LOF are outlier
    detectors, so a literal reading of ``at least two types (classification, clustering,
    regression)'' may not credit them. Decision D-A2-1 fixes this.
  \item \textbf{Labelled data.} With no labels, the brief's accuracy, precision, recall and F1
    cannot be reported honestly. Decision D-A2-3 fixes this.
  \item \textbf{More datasets.} A1 used one source. The brief says one simple dataset earns ``a
    basic score''.
  \item \textbf{Any GPU run of the language models}, and any transaction-sequence code at all.
\end{enumerate}

\subsection{What the unit's own material says to do}

The weekly material points in the same direction as this plan, and the report can say so:
\begin{itemize}
  \item \textbf{Week 4 workshop.} Clean, anonymise sensitive fields, engineer security features,
    handle imbalance with class weights, and prefer recall-oriented metrics. Its closing
    self-review asks for a \emph{data card} and a \emph{model card}; both are in
    Section~\ref{sec:packaging}.
  \item \textbf{Week 5.} Extract, transform, load; normalise and standardise; join and
    concatenate; visualise class balance, distributions, correlations and time series. Its
    licence slide says HD students should use additional datasets and respect their licences.
  \item \textbf{Week 6.} Regression is evaluated with MSE, RMSE, MAE and $R^2$; K-Means with
    inertia, silhouette and Davies--Bouldin. The workshop calls outliers ``not always attacks but
    worth investigation''. It gives the workflow ``cluster → sample validation → derive patterns
    → propose simple rules''. This is the project's framing in the lecturer's own words.
  \item \textbf{Week 7.} DBSCAN with $\text{MinPts}=2\times D$ and a k-distance graph for
    $\varepsilon$, with noise points read as anomalies. Logistic regression and KNN on
    standardised data, scaled with the training set's parameters only. Accuracy, precision,
    recall, F1 and the confusion matrix. Decision trees and random forests appear as
    ``additional'' models.
\end{itemize}
